\section{Live Tool Timing, Latency and Complete Tunnel Controls (E13)}\label{sec:r6-supp}
\subsection{Questions and measured boundary}
E13 executes the saved 690 real-agent action trajectories again through a live TLS tool server. It asks whether arrival timing restores intent inference after length/count shaping, whether an observer can identify the end flush or cover requests, and what actual waiting time accompanies each intervention. Six conditions yield 4,140 newly measured sessions. These are new tool executions and ciphertext arrival measurements, not new LLM generations or a new customer population.

The client driver waits for each real response before releasing its next saved action. The server invokes the pinned retail tool implementation against a live database clone; cover invokes the same implementation against a separate clone. Each real response is checked against the corresponding saved model observation, and the final live database is compared with unprotected execution of the same actions. This feedback-synchronized tool loop measures real queues and service times while holding the planner's action sequence fixed. Its clocks exclude LLM/user generation and WAN transit. No artificial read/write delay is injected. Reward preservation is inherited when all responses and final live state agree; it is not a new stochastic planner trial.

\subsection{Transport and execution controls}
All conditions use a fresh TLS~1.2 connection with ECDHE-RSA-AES128-GCM-SHA256 and TCP\_NODELAY through a byte-transparent loopback relay. The relay timestamps ciphertext chunk arrivals using a monotonic high-resolution clock and parses records; a record inherits the timestamp of its final observed chunk. The retained chunk log permits inspection of coalescing. These are relay-observed record times, not kernel packet timestamps. The application projection excludes handshakes and closure alerts. One session is measured at a time; episode and condition order are shuffled using a fixed seed.

The six conditions are raw transport, per-call constant framing, serial shadow cover with constant framing, the same cover with at most two outstanding RPCs, and two complete fixed-rate tunnel settings. The cover scheduler uses E9's frozen training donors, count template and random seed, observes only the executed prefix, and flushes remaining counts after the last real response. Real and shadow states have separate execution locks; neither a marker nor a shadow response is exposed to the passive link observer.

The tunnel fragments and reassembles the complete application stream. Both directions send a 4,096-byte plaintext cell at every slot, including when no payload is ready. An encrypted four-byte field gives the useful payload length; the remainder is padding. Dummy cells never invoke a tool. Every cell produces a 4,120-byte TLS ciphertext record, plus its five-byte header. This implements the fixed-size/constant-rate traffic-flow principle described by RFC~9347~\cite{rfc9347} over TLS; it is not an IPsec wire-format implementation.

The common service horizon is fixed using training messages only. For each training trajectory, sum the required request/response payload cells plus four slack cells per real call. The maximum budget, multiplied by 1.5 and rounded up with 16 additional slots, is lower-bounded by 128. This yields 133 slots: 266\,ms at 2\,ms per cell or 665\,ms at 5\,ms. Both directions continue through this horizon after the real workload completes. A horizon overflow is a retained failure, never a secret-dependent extension. The per-call framing caps use E9's frozen caps plus 256 bytes for the RPC envelope, with whole-cap overflow buckets. The operating points are fixed before test measurement.

Actual deadlines are rechecked against the measurement clock after sleeping; no ideal timestamps replace late sends or catch-up bursts. Both receivers must consume all 133 cells before closure. The saved outputs include cell deadlines, observed sends, pending queues and all requests/completions. A host scheduler and a local Python gateway remain part of this measured implementation. The two tunnel settings are fixed operating points, not an optimized cost frontier or a measurement of multi-session amortization.

\subsection{Attacks and fitting order}
The original customer-disjoint split is preserved: 414 training trajectories from 29 customers and 276 test trajectories from 24 customers. All 2,484 final training measurements precede attack fitting and a model-hash freeze; only then are the 1,656 final test sessions collected. An earlier complete development pass exposed two incomplete tunnel tails at connection closure. We added a receive-drain barrier, restored E9-style profilers to the attacker suite, and recollected the entire cohort; no individual test sessions were selectively replaced. Development records are archived separately. The original trajectories and labels were already known from E9, so this is a post hoc measurement study on the same population.

The intent suite has a training-majority baseline, balanced logistic regression ($C=1$), gradient boosting with 100 trees, and random forests with 200 trees and a minimum leaf size of two. Its three basic views are size/order, timing/direction and their joint features. Features include direction-specific count, size/gap quantiles, request/response size histograms, exchange-length vectors, duration, gap thresholds at 0.1/1/5/20\,ms, the first/last 32 positions, and largest-gap change-point summaries. For non-tunnel conditions, training-only size-to-tool profilers additionally feed a decoded gradient-boosting attacker and a training-selected count rule; these join both size/order and full-suite evaluations. At test time, all attackers consume observable records only. True tool names, RPC IDs, real/cover markers, true flush boundaries and internal execution times are excluded from test features.

Separate 80-tree forests (depth eight, minimum leaf size five) detect the flush suffix and real versus cover requests. Flush detection uses observable record neighborhoods and relative position/time; cover detection uses the same visible features at request-record positions. True execution annotations supply fitting/scoring labels only. For the flush-stripped intent view, training traces are transformed by three-fold customer-disjoint out-of-fold detectors. The full-training detector transforms test traces; a maximum-likelihood single transition estimates the cut. No oracle test boundary is used. The full adaptive suite unions every applicable view, allowing an observer to ignore timing or retain the unstripped trace.

\subsection{Inference, latency and reproducibility}
Intent intervals use 2,000 common customer-cluster bootstrap draws, retaining all tasks and repeats and recomputing the maximum over the suite in each draw. Paired differences reuse these draws. Draws lacking either target class have undefined BA/AUROC and are omitted from interval quantiles; their counts are archived. Intervals condition on fitted predictions and do not supply a new independent population estimate. The E9 absolute criterion remains unchanged: BA upper 95\% bound below 0.60 and AUROC below 0.65. Detector BA also uses customer clusters; detector AUROC is reported as a point estimate.

Latency is reported for individual real RPCs, the first real response, the last real response and the complete traffic window. Client scheduling queues, backend service, slot lateness and timeouts are retained separately. The descriptive count of real calls exceeding 100\,ms is a reference threshold, not a service-level guarantee. Byte ratios count measured application TLS records and their headers; call ratios count actual tool RPCs, excluding tunnel dummy cells. Source, calibration, model and capture hashes bind the analysis. The saved captures, fitted models and predictions support offline verification without new LLM or network requests.

\subsection{Findings across the measured boundaries}
For baseline prompts, serial cover's full-suite cancellation BA exceeds its size/order BA by 11.4 points [3.7,19.6]. The full-suite reduction relative to raw remains 30.3 points [23.1,37.6] for cancellation and 32.9 [19.9,40.7] for return. Thus the measurement supports a substantial reduction while exposing a residual clock signal. Flush and real/cover detectors pool both prompt arms: serial/parallel flush BA is 90.4\%/92.3\%, and cover BA is 82.5\%/96.1\%. Successful segmentation does not itself establish successful intent recovery; the separately fitted stripped-trace attacks test that transfer.

All 1,380 tunnel sessions contain exactly 133 records per direction with identical ciphertext lengths. Nevertheless, the measured directions can interleave differently, and scheduling is imperfect. In baseline-prompt tests, 2/5\,ms tunnel slot lateness has p95 1.88/1.68\,ms and maxima 282.6/307.5\,ms. These retained observations characterize the implementation, rather than an ideal zero-jitter clock. The full-suite baseline BAs are 62.5\%/52.9\% and 53.1\%/57.3\%; neither setting meets the two-target absolute criterion. Their 125.7$\times$ byte cost and one tool call per original action identify a different resource trade-off from shadow cover's 15.1$\times$ bytes and 3.83$\times$ calls. These fixed operating points do not establish an optimal tunnel or universal superiority of either mechanism.

No formal session fails, loses a real response, changes final live state or sends a cell before its deadline. Twelve of 2,000 cancellation bootstrap draws do not contain both target classes; all return draws are defined. The archived invalid-draw counters distinguish this finite-sample calculation issue from execution failures.

\subsection{Evaluation under WAN transit delay and jitter (Exp 4 Extension)}\label{sec:wan-netem-eval}
Local loopback measurement isolates backend execution and tool scheduling from transit variance, but real agent deployments operate over wide-area networks subject to packet propagation delays and jitter. To evaluate whether timing-based flush detection and intent inference hold outside the loopback environment, we extend Exp~4 with network emulation matching standard wide-area link profiles:
\begin{itemize}
    \item \textbf{WAN Mild}: 20\,ms one-way delay with 5\,ms normal jitter ($\sigma = 5$\,ms), representing regional cloud tool hosting ($40$\,ms nominal RTT).
    \item \textbf{WAN Realistic}: 50\,ms one-way delay with 10\,ms normal jitter ($\sigma = 10$\,ms), representing cross-region or transcontinental tool hosting ($100$\,ms nominal RTT).
\end{itemize}
The emulation runs directly in the bidirectional TLS transport relay using high-precision non-reordering FIFO delivery queues matching the behavior of Linux \texttt{tc netem}. Each real tool response across the WAN link gates the next action release, faithfully preserving feedback synchronization against live database clones. We repeat measurements on the complete set of 276 customer-held-out test episodes across both prompt arms (standard baseline and system-prompt rules), yielding 1,104 new TLS sessions per WAN profile (2,208 additional measured sessions total) across raw transport, per-call constant framing, serial shadow cover, and 2\,ms complete tunnels.

\paragraph{Resilience of flush detection}
Table~\ref{tab:wan-flush-comparison} assesses flush-phase detection across 15,278 ciphertext records from serial shadow cover sessions. We evaluate two attacker postures: (1) \emph{Direct Transfer}, where the attacker applies the frozen loopback detector $D_{\text{loop}}$ without network-specific recalibration, and (2) \emph{Adapted 5-Fold CV}, where the attacker refits the 80-tree random forest on customer-grouped training folds under the observed WAN network condition.

Under both WAN conditions, flush detection exhibits remarkable resilience. With direct transfer of $D_{\text{loop}}$, the detector achieves 89.4\% point BA (95\% CI [87.2\%, 91.7\%], AUROC 0.962) under 20\,ms/5\,ms jitter and 89.4\% point BA (95\% CI [87.3\%, 91.7\%], AUROC 0.963) under 50\,ms/10\,ms jitter, falling by only 1.0 percentage point from the loopback baseline (90.4\% BA, AUROC 0.967). Adapted retraining yields consistent performance (89.2\% and 89.0\% BA). This stability occurs because the flush phase transitions from feedback-gated tool intervals (punctuated by real database queries and network RTTs) to an un-gated burst of dummy requests. This macroscopic transition in inter-request density and duration creates a structural signature that easily cuts through 5--10\,ms transit noise.
\begin{table}[t]
\centering
\caption{Serial shadow cover flush-phase detection across network conditions (15,278 records from 276 test episodes). Direct transfer applies the frozen loopback detector $D_{\text{loop}}$; Adapted fits 5-fold customer-grouped models on the WAN workload.}
\label{tab:wan-flush-comparison}
\small
\begin{tabular}{lcccc}
\toprule
\textbf{Network Condition} & \textbf{Model} & \textbf{Flush Point BA} & \textbf{95\% CI} & \textbf{AUROC} \\
\midrule
Loopback (0 ms, 0 ms jitter) & Frozen $D_{\text{loop}}$ & 90.4\% & [88.3, 92.5] & 0.967 \\
\midrule
WAN Mild (20 ms, 5 ms jitter) & Transfer $D_{\text{loop}}$ & 89.4\% & [87.2, 91.7] & 0.962 \\
& Adapted $D_{\text{WAN}}$ & 89.2\% & [86.6, 91.5] & 0.958 \\
\midrule
WAN Realistic (50 ms, 10 ms jitter) & Transfer $D_{\text{loop}}$ & 89.4\% & [87.3, 91.7] & 0.963 \\
& Adapted $D_{\text{WAN}}$ & 89.0\% & [86.5, 91.3] & 0.958 \\
\bottomrule
\end{tabular}
\end{table}

\paragraph{Intent classification under transit jitter}
Table~\ref{tab:wan-controls-comparison} compares full-suite intent Balanced Accuracy across loopback and WAN conditions. For raw transport and per-call constant framing, intent classification is dominated by packet size and sequence order distributions, which are network-delay invariant; hence cancellation BA remains 91.7\% and 68.6\%, respectively.

For serial shadow cover, however, the timing signal exploited on loopback is substantially attenuated. On loopback, residual execution time differences allowed the full attack union to elevate cancellation BA to 61.4\% (11.4 points above size/order). Under WAN conditions, random transit jitter washes out microsecond-level tool service variations: cancellation BA falls to 50.0\% (exact chance level) under 20\,ms/5\,ms jitter and 54.0\% under 50\,ms/10\,ms jitter; return BA remains 50.7\% and 54.6\%. Thus, WAN network jitter acts as natural timing obfuscation against micro-timing intent inference, strengthening the intent-hiding property of serial cover while leaving macroscopic flush boundaries visible.

\paragraph{Operational limits: tunnel horizon exhaustion}
The WAN extension reveals a critical operational vulnerability in fixed-horizon tunnel designs. The 2\,ms full tunnel enforces a strict 133-slot horizon (266\,ms) sized against local loopback execution. On loopback, all 276 sessions successfully completed all tool calls within this window.

When executed over WAN links, each sequential tool RPC incurs network RTT. Under 20\,ms one-way delay (40\,ms RTT per call), tasks requiring multiple tool round-trips exhaust the 133 slots before their responses arrive; 191 out of 276 sessions (69.2\%) fail to complete due to horizon exhaustion (only 85/276 complete). Under 50\,ms delay (100\,ms RTT), horizon exhaustion surges to 98.2\% (271 out of 276 sessions fail; only 5 single-call sessions complete). This highlights an intrinsic deployment hazard: fixed-rate constant-bandwidth tunnels require worst-case horizon provisioning over WAN links, which would drastically multiply byte overhead (e.g., thousands of slots), or else suffer catastrophic protocol failure whenever network latency fluctuates.

In contrast, serial shadow cover and per-call constant framing adapt naturally to network delays: all 276 sessions complete successfully with 100\% response and database state equality across both WAN conditions. Their real-chain p95 latency expands predictably with network round-trips: from 26.2\,ms on loopback to 855.4\,ms under 20\,ms WAN and 1999.5\,ms under 50\,ms WAN.
\begin{table*}[t]
\centering
\caption{Timing defenses under loopback versus WAN conditions (276 test episodes). Full-suite intent Balanced Accuracy (BA) reported for cancellation / return. Incomplete/exhausted reflects sessions failing to complete all real calls within the fixed horizon.}
\label{tab:wan-controls-comparison}
\footnotesize\setlength{\tabcolsep}{3pt}
\begin{tabular}{lcccccc}
\toprule
\textbf{Condition} & \textbf{Network Environment} & \textbf{Full-Suite BA} & \textbf{Real p95 ms} & \textbf{Dur p95 ms} & \textbf{Byte Mult.} & \textbf{Status / Exh.} \\
\midrule
Raw transport & Loopback (0 ms, 0 ms jit.) & 91.7\% / 86.4\% & 15.4 & 15.5 & 1.0$\times$ & 276/276 \\
& WAN Mild (20 ms, 5 ms jit.) & 91.7\% / 86.4\% & 506.4 & 506.4 & 1.0$\times$ & 276/276 \\
& WAN Real. (50 ms, 10 ms jit.) & 91.7\% / 86.4\% & 1178.8 & 1178.8 & 1.0$\times$ & 276/276 \\
\midrule
Per-call constant framing & Loopback (0 ms, 0 ms jit.) & 68.6\% / 62.7\% & 16.6 & 16.7 & 3.9$\times$ & 276/276 \\
& WAN Mild (20 ms, 5 ms jit.) & 68.6\% / 54.8\% & 504.8 & 504.9 & 3.9$\times$ & 276/276 \\
& WAN Real. (50 ms, 10 ms jit.) & 68.6\% / 54.8\% & 1172.6 & 1172.6 & 3.9$\times$ & 276/276 \\
\midrule
Serial shadow cover & Loopback (0 ms, 0 ms jit.) & 61.4\% / 53.4\% & 26.2 & 46.6 & 15.1$\times$ & 276/276 \\
& WAN Mild (20 ms, 5 ms jit.) & 50.0\% / 50.7\% & 855.4 & 1314.9 & 15.3$\times$ & 276/276 \\
& WAN Real. (50 ms, 10 ms jit.) & 54.0\% / 54.6\% & 1999.5 & 3112.4 & 15.3$\times$ & 276/276 \\
\midrule
Full tunnel (2 ms cell) & Loopback (0 ms, 0 ms jit.) & 62.5\% / 52.9\% & 56.4 & 270.4 & 125.7$\times$ & 276/276 \\
& WAN Mild (20 ms, 5 ms jit.) & 54.4\% / 56.8\% & 286.7 & 291.6 & 127.0$\times$ & 85/276 \\
& WAN Real. (50 ms, 10 ms jit.) & 55.0\% / 52.6\% & 269.5 & 325.9 & 127.0$\times$ & 5/276 \\
\bottomrule
\end{tabular}
\end{table*}

\clearpage
\begin{table}[htbp]\centering\footnotesize
\caption{E13 baseline prompt: intent BA [95\% customer-cluster interval] and AUROC. The full suite unions all available views.}
\begin{tabular}{llrrrr}\toprule
Condition & View & Cancel BA [CI] & Return BA [CI] & Cancel AUC & Return AUC\\\midrule
Raw & Size/order & \mbox{91.7 [84.7, 100.0]} & \mbox{86.4 [79.5, 97.2]} & 0.917 & 0.875 \\
Raw & Timing/direction & \mbox{53.8 [50.0, 72.5]} & \mbox{52.9 [50.0, 64.1]} & 0.531 & 0.570 \\
Raw & Joint features & \mbox{84.3 [72.5, 97.4]} & \mbox{80.5 [69.9, 93.6]} & 0.836 & 0.835 \\
Raw & Full adaptive suite & \mbox{91.7 [84.7, 100.0]} & \mbox{86.4 [79.5, 97.2]} & 0.917 & 0.875 \\
\addlinespace[2pt]
Per-call padding & Size/order & \mbox{68.6 [50.0, 92.8]} & \mbox{54.8 [50.0, 74.8]} & 0.678 & 0.514 \\
Per-call padding & Timing/direction & \mbox{54.6 [50.0, 63.0]} & \mbox{62.7 [53.8, 78.0]} & 0.500 & 0.657 \\
Per-call padding & Joint features & \mbox{51.0 [50.0, 61.4]} & \mbox{50.2 [50.0, 69.0]} & 0.569 & 0.545 \\
Per-call padding & Full adaptive suite & \mbox{68.6 [50.8, 92.8]} & \mbox{62.7 [54.2, 79.3]} & 0.678 & 0.657 \\
\addlinespace[2pt]
Shadow, serial & Size/order & \mbox{50.0 [50.0, 56.9]} & \mbox{50.7 [50.0, 62.5]} & 0.500 & 0.505 \\
Shadow, serial & Timing/direction & \mbox{61.4 [53.8, 67.7]} & \mbox{52.9 [50.0, 62.5]} & 0.625 & 0.500 \\
Shadow, serial & Joint features & \mbox{51.0 [50.0, 64.6]} & \mbox{50.0 [50.0, 57.7]} & 0.596 & 0.500 \\
Shadow, serial & Estimated flush removed & \mbox{57.8 [50.0, 70.7]} & \mbox{53.4 [50.0, 64.6]} & 0.583 & 0.530 \\
Shadow, serial & Full adaptive suite & \mbox{61.4 [54.6, 71.1]} & \mbox{53.4 [50.0, 65.8]} & 0.625 & 0.530 \\
\addlinespace[2pt]
Shadow, two in flight & Size/order & \mbox{51.7 [50.0, 68.3]} & \mbox{57.7 [51.5, 66.1]} & 0.509 & 0.569 \\
Shadow, two in flight & Timing/direction & \mbox{50.0 [50.0, 57.1]} & \mbox{50.0 [50.0, 59.9]} & 0.554 & 0.500 \\
Shadow, two in flight & Joint features & \mbox{50.0 [50.0, 57.1]} & \mbox{52.6 [50.0, 63.3]} & 0.531 & 0.522 \\
Shadow, two in flight & Estimated flush removed & \mbox{50.0 [50.0, 64.0]} & \mbox{50.0 [50.0, 56.7]} & 0.500 & 0.500 \\
Shadow, two in flight & Full adaptive suite & \mbox{51.7 [50.0, 69.1]} & \mbox{57.7 [53.4, 66.8]} & 0.554 & 0.569 \\
\addlinespace[2pt]
Full tunnel, 2 ms & Size/order & \mbox{62.5 [54.3, 70.1]} & \mbox{52.9 [50.0, 61.0]} & 0.594 & 0.521 \\
Full tunnel, 2 ms & Timing/direction & \mbox{53.1 [50.2, 59.3]} & \mbox{50.0 [50.0, 55.6]} & 0.641 & 0.500 \\
Full tunnel, 2 ms & Joint features & \mbox{56.2 [50.9, 63.9]} & \mbox{50.0 [50.0, 57.4]} & 0.656 & 0.569 \\
Full tunnel, 2 ms & Full adaptive suite & \mbox{62.5 [55.6, 70.2]} & \mbox{52.9 [50.1, 61.3]} & 0.656 & 0.569 \\
\addlinespace[2pt]
Full tunnel, 5 ms & Size/order & \mbox{51.4 [50.0, 59.0]} & \mbox{54.5 [50.0, 60.3]} & 0.500 & 0.513 \\
Full tunnel, 5 ms & Timing/direction & \mbox{53.1 [50.0, 71.9]} & \mbox{53.2 [50.0, 62.4]} & 0.575 & 0.548 \\
Full tunnel, 5 ms & Joint features & \mbox{52.4 [50.0, 61.8]} & \mbox{57.3 [53.1, 63.7]} & 0.534 & 0.555 \\
Full tunnel, 5 ms & Full adaptive suite & \mbox{53.1 [50.3, 71.9]} & \mbox{57.3 [54.2, 64.2]} & 0.575 & 0.555 \\
\addlinespace[2pt]
\bottomrule\end{tabular}
\end{table}
\clearpage
\begin{table}[htbp]\centering\footnotesize
\caption{E13 rules prompt: intent BA [95\% customer-cluster interval] and AUROC. The full suite unions all available views.}
\begin{tabular}{llrrrr}\toprule
Condition & View & Cancel BA [CI] & Return BA [CI] & Cancel AUC & Return AUC\\\midrule
Raw & Size/order & \mbox{94.4 [83.3, 100.0]} & \mbox{87.0 [78.0, 97.6]} & 0.944 & 0.905 \\
Raw & Timing/direction & \mbox{57.8 [50.0, 70.6]} & \mbox{56.1 [50.0, 63.9]} & 0.598 & 0.557 \\
Raw & Joint features & \mbox{84.7 [75.2, 98.2]} & \mbox{79.1 [67.8, 92.0]} & 0.851 & 0.822 \\
Raw & Full adaptive suite & \mbox{94.4 [83.3, 100.0]} & \mbox{87.0 [78.0, 97.6]} & 0.944 & 0.905 \\
\addlinespace[2pt]
Per-call padding & Size/order & \mbox{70.1 [50.0, 86.0]} & \mbox{54.8 [50.0, 74.7]} & 0.680 & 0.500 \\
Per-call padding & Timing/direction & \mbox{55.8 [50.0, 70.2]} & \mbox{50.6 [50.0, 62.4]} & 0.530 & 0.534 \\
Per-call padding & Joint features & \mbox{59.4 [50.0, 76.6]} & \mbox{53.6 [50.0, 71.0]} & 0.575 & 0.508 \\
Per-call padding & Full adaptive suite & \mbox{70.1 [50.0, 86.0]} & \mbox{54.8 [50.0, 74.8]} & 0.680 & 0.534 \\
\addlinespace[2pt]
Shadow, serial & Size/order & \mbox{58.2 [50.2, 71.2]} & \mbox{57.9 [50.7, 66.8]} & 0.581 & 0.597 \\
Shadow, serial & Timing/direction & \mbox{51.5 [50.0, 69.8]} & \mbox{57.7 [52.3, 68.8]} & 0.500 & 0.560 \\
Shadow, serial & Joint features & \mbox{51.1 [50.0, 64.4]} & \mbox{55.1 [52.2, 65.3]} & 0.519 & 0.555 \\
Shadow, serial & Estimated flush removed & \mbox{50.1 [50.0, 55.6]} & \mbox{60.7 [50.8, 70.9]} & 0.500 & 0.622 \\
Shadow, serial & Full adaptive suite & \mbox{58.2 [52.9, 72.2]} & \mbox{60.7 [56.3, 71.7]} & 0.581 & 0.622 \\
\addlinespace[2pt]
Shadow, two in flight & Size/order & \mbox{54.7 [50.0, 66.7]} & \mbox{60.2 [50.5, 73.8]} & 0.580 & 0.576 \\
Shadow, two in flight & Timing/direction & \mbox{50.0 [50.0, 57.8]} & \mbox{54.0 [50.0, 62.6]} & 0.556 & 0.585 \\
Shadow, two in flight & Joint features & \mbox{55.3 [50.0, 67.1]} & \mbox{52.3 [50.0, 61.0]} & 0.578 & 0.546 \\
Shadow, two in flight & Estimated flush removed & \mbox{50.8 [50.0, 61.7]} & \mbox{54.0 [50.0, 67.0]} & 0.541 & 0.528 \\
Shadow, two in flight & Full adaptive suite & \mbox{55.3 [50.8, 68.2]} & \mbox{60.2 [54.5, 73.9]} & 0.580 & 0.585 \\
\addlinespace[2pt]
Full tunnel, 2 ms & Size/order & \mbox{51.5 [50.0, 65.0]} & \mbox{53.7 [50.0, 59.7]} & 0.500 & 0.500 \\
Full tunnel, 2 ms & Timing/direction & \mbox{53.1 [50.0, 64.2]} & \mbox{53.1 [50.0, 60.9]} & 0.500 & 0.563 \\
Full tunnel, 2 ms & Joint features & \mbox{51.5 [50.0, 58.1]} & \mbox{59.1 [52.2, 66.1]} & 0.500 & 0.590 \\
Full tunnel, 2 ms & Full adaptive suite & \mbox{53.1 [50.0, 65.6]} & \mbox{59.1 [52.8, 66.1]} & 0.500 & 0.590 \\
\addlinespace[2pt]
Full tunnel, 5 ms & Size/order & \mbox{50.0 [50.0, 53.6]} & \mbox{55.4 [50.9, 62.7]} & 0.500 & 0.549 \\
Full tunnel, 5 ms & Timing/direction & \mbox{50.0 [50.0, 50.4]} & \mbox{56.3 [51.0, 61.6]} & 0.500 & 0.530 \\
Full tunnel, 5 ms & Joint features & \mbox{50.0 [50.0, 57.5]} & \mbox{52.3 [50.0, 59.5]} & 0.500 & 0.534 \\
Full tunnel, 5 ms & Full adaptive suite & \mbox{50.0 [50.0, 57.5]} & \mbox{56.3 [55.0, 63.3]} & 0.500 & 0.549 \\
\addlinespace[2pt]
\bottomrule\end{tabular}
\end{table}
\clearpage
\begin{table}[htbp]\centering\footnotesize
\caption{E13 live execution costs, both prompt arms. Medians and 95th percentiles describe observed local RPCs; they are not production-service latency estimates.}
\begin{tabular}{llrrrrrr}\toprule
Prompt & Condition & Bytes$\times$ & Calls$\times$ & First p50/p95 & Last real p50/p95 & Traffic p50/p95 & RPC p50/p95\\\midrule
baseline & Raw & 1.00 & 1.00 & 2.3/3.4 & 8.8/15.4 & 8.8/15.5 & 1.0/1.8\\
baseline & Per-call padding & 3.92 & 1.00 & 2.3/3.7 & 8.5/16.6 & 8.5/16.7 & 1.0/1.8\\
baseline & Shadow, serial & 15.11 & 3.83 & 3.0/5.0 & 12.5/26.2 & 27.8/46.6 & 0.9/1.7\\
baseline & Shadow, two in flight & 15.11 & 3.83 & 2.7/4.5 & 11.3/20.8 & 23.6/37.4 & 1.2/2.4\\
baseline & Full tunnel, 2 ms & 125.71 & 1.00 & 5.4/11.2 & 30.9/56.4 & 267.0/270.4 & 4.0/6.2\\
baseline & Full tunnel, 5 ms & 125.71 & 1.00 & 11.3/12.4 & 71.4/112.3 & 666.1/666.6 & 10.0/10.9\\
rules & Raw & 1.00 & 1.00 & 2.2/3.7 & 8.1/15.7 & 8.2/15.7 & 1.0/1.7\\
rules & Per-call padding & 3.95 & 1.00 & 2.2/3.6 & 8.4/15.5 & 8.4/15.5 & 1.0/1.8\\
rules & Shadow, serial & 15.47 & 3.88 & 3.0/5.2 & 12.7/23.9 & 28.9/41.8 & 1.0/1.7\\
rules & Shadow, two in flight & 15.47 & 3.88 & 2.7/4.4 & 11.1/20.7 & 23.8/36.5 & 1.2/2.4\\
rules & Full tunnel, 2 ms & 128.38 & 1.00 & 5.4/7.7 & 29.8/48.8 & 267.0/267.9 & 4.0/5.8\\
rules & Full tunnel, 5 ms & 128.38 & 1.00 & 11.1/12.4 & 69.5/112.9 & 666.2/666.6 & 10.0/11.0\\
\bottomrule\end{tabular}
\par\smallskip All time columns are milliseconds. The saved results additionally report queue/service times, a descriptive 100 ms per-call deadline count, cell lateness, failures, and paired mean last-response differences.
\end{table}
\begin{table}[htbp]\centering\footnotesize
\caption{E13 learned timing/shape detectors on held-out customers, pooling both prompt arms. Flush labels identify records after the online flush declaration; cover labels identify encrypted requests. Ground truth is used only for fitting/scoring, never as an attacker feature.}
\begin{tabular}{llrrr}\toprule
Condition & Detector & BA [95\% CI] & AUROC & Records/requests\\\midrule
Shadow, serial & flush & \mbox{90.4 [88.3, 92.5]} & 0.967 & 15278 \\
Shadow, serial & cover & \mbox{82.5 [80.3, 84.7]} & 0.898 & 7639 \\
Shadow, two in flight & flush & \mbox{92.3 [90.6, 93.9]} & 0.979 & 15278 \\
Shadow, two in flight & cover & \mbox{96.1 [95.4, 96.7]} & 0.995 & 7639 \\
\bottomrule\end{tabular}
\end{table}

